\documentclass[10pt,a4paper]{amsart}
\usepackage[margin=19mm]{geometry}
\usepackage{amsmath,amssymb,mathtools}
\usepackage[scr=rsfs,cal=euler]{mathalfa}
\usepackage{xfrac}
\usepackage[unicode,hidelinks,bookmarks=false]{hyperref}
\newcommand{\ie}{i.e.\ }
\newcommand{\cf}{cf.\ }
\newcommand{\resp}{respectively\ }
\newcommand{\Subsection}{\subsection}
\providecommand{\nobreakdash}{\nobreak\hbox{-}}
\setlength{\emergencystretch}{2em}
\multlinegap0pt
\usepackage{bm}
\usepackage{bbm}
\usepackage{dsfont}
\usepackage{xspace}
\usepackage{cancel}
\usepackage{mathtools}
\usepackage{amsthm}
\makeatletter
\@ifundefined{definition}{\newtheorem{definition}{Definition}}{}
\@ifundefined{lemma}{\newtheorem{lemma}[definition]{Lemma}}{}
\makeatother
\makeatletter
\newcommand{\gc}{i_{g_{\frac{1}{c}}}}
\expandafter\ifx\csname tehtratk\endcsname\relax
\newenvironment{tehtratk}%
{\begin{list}{\arabic{enumi}.}{\usecounter{enumi}%
			\setlength{\labelsep}{0.5em}%
			\settowidth{\labelwidth}{(\arabic{enumi})}%
			\setlength{\leftmargin}{\labelwidth+\labelsep}}}%
	{\end{list}}
\fi
\makeatother
\title[The Momentum Light Ray Transform]{The Momentum Light Ray Transform}
\author{Sombuddha Bhattacharyya, Tuhin Mondal, Suman Kumar Sahoo}
\date{Source: arXiv:2510.18450v1 (CC BY 4.0)}
\begin{document}
\maketitle

\noindent\emph{Source and licence.} The Momentum Light Ray Transform. Version arXiv:2510.18450v1 (CC BY 4.0), distributed under the Creative Commons Attribution 4.0 International licence, as recorded in the arXiv licence metadata for this exact version and shown on the abstract page https://arxiv.org/abs/2510.18450v1. Full rights notice: licenses/arxiv-2510.18450-CC-BY-4.0.txt.\par\smallskip

\noindent\textit{Editorial scope.} Counted unit: the Lemma whose source label is \texttt{lem\_commutator} in the article (source0.tex lines 282--289, source label \texttt{lem\_commutator}), delivered together with the article's own complete original proof of it (source0.tex lines 290--383, one \texttt{proof} environment of 94 lines). No mathematics is written, repaired or re-derived: statement and proof are the article's own text line for line. Required context retained uncounted: none. Every definition, display and label of those lines is retained unchanged. Omitted from this package and not redistributed: 1--281, 384--1698. Nothing inside a retained excerpt is omitted or altered: every retained body is delivered exactly as the article's lines are written, so no omission receipt of any kind accompanies this package. Citations: the retained excerpts carry no citation command at all, so no bibliography entry of the article is transcribed and no reference is redistributed. Reference adaptation: none was needed and none was made; every reference inside the delivered unit resolves inside it, so no unresolved reference or citation is produced. Printed slips kept literal: no printed word, symbol, hypothesis or number of the counted statement or of its proof was altered. Layout and class: the article's own class is \texttt{amsart} with packages including amsmath, amscd, mathtools, amssymb, amsthm, enumerate, amsrefs, hyperref, graphicx, mathrsfs, nccmath, cancel, calc, comment; the wrapper is the lane's pinned \texttt{amsart} editorial wrapper at 10pt on A4, which restates the theorem-like environments the retained text needs and adds the article's own macro definitions that the retained text needs (\texttt{\textbackslash gc}), with extra wrapper packages (bm, bbm, dsfont, xspace, cancel, mathtools). The delivered numbering is the unit's own. Assets: no retained span contains a figure environment, a graphics-inclusion command, a PDF-page inclusion, a table or a TikZ picture; no image file is copied into this package, the package declares no asset dependency and no image file is redistributed with it. Licence: version arXiv:2510.18450v1 of this article is distributed under the Creative Commons Attribution 4.0 International licence, as recorded in the arXiv licence metadata for this exact version and shown on the abstract page https://arxiv.org/abs/2510.18450v1. A full rights notice is delivered at licenses/arxiv-2510.18450-CC-BY-4.0.txt. Characters: every retained excerpt is pure ASCII, so the delivered engine, XeLaTeX with its default Latin Modern text font, reports no missing character.
	\begin{lemma}\label{lem_commutator}
Let $n\ge 2$ and   $f^{(m-2)} \in S^{m-2}$ for $m\ge 2$, then the following relations hold.
\begin{itemize}
\item [1.]	\[J\gc f^{(m-2)}= D_{m}\gc Jf^{(m-2)}+ C_mf^{(m-2)},\] 
where $D_m=\binom{m}{2}^{-1}\binom{m-2}{2},$ and $C_m=\binom{m}{2}^{-1}(n+m-1)$.
\item [2.] \[\gc f^{(m-2)}=0 \implies  f^{(m-2)}=0.\]
\end{itemize}
	\end{lemma}

	\begin{proof}
		 For fixed indices $i_1,i_2,\dots,i_{m-2} = 0,\cdots,n$,
		\begin{equation*}
			\begin{aligned}
				\Big[J\gc f^{(m-2)} \Big]_{i_1\dots i_{m-2}}&=-c^2\Big[\gc f^{(m-2)} \Big]_{i_1\dots i_{m-2}00}+\sum_{k=1}^n\Big[\gc f^{(m-2)} \Big]_{i_1\dots i_{m-2}kk}\\
				&=-c^2\binom{m}{2}^{-1}\underbrace{\Big[g_{i_1i_2}f^{(m-2)}_{i_3\dots i_{m-2}00}+\cdots +g_{i_{m-1}i_m}f^{(m-2)}_{i_1\dots i_{m-2}} \Big]}_{i_{m-1}=i_m=0}\\
				&\quad+\sum_{k=1}^{n}\binom{m}{2}^{-1}\underbrace{\Big[g_{i_1i_2}f^{(m-2)}_{i_3\dots i_{m-2}kk}+\cdots +g_{i_{m-1}i_m}f^{(m-2)}_{i_1\dots i_{m-2}} \Big]}_{i_{m-1}=i_m=k}.
			\end{aligned}
		\end{equation*}
		
The last term in the first bracket is equal to 
\[
-c^2 \binom{m}{2}^{-1} g_{00} f^{(m-2)}_{i_1 \dots i_{m-2}}
= \binom{m}{2}^{-1} f^{(m-2)}_{i_1 \dots i_{m-2}}.
\]
A similar calculation shows that the last term in each of the other brackets is also equal to 
		\(\binom{m}{2}^{-1} f^{(m-2)}_{i_1 \dots i_{m-2}}\). 
		Therefore, these terms contribute a total of 
		\((n+1)\binom{m}{2}^{-1} f^{(m-2)}\).
		Next, we carefully handle the remaining terms in each bracket.  
		First, consider those terms where \(i_{m-1} = i_m = 0\), so that they remain with \(f^{(m-2)}\). 
		In the first bracket, their contribution is 
		\[
		-c^2 \binom{m}{2}^{-1} \Big[
		g_{i_1 i_2}(f^{(m-2)}_{00})_{i_3 \dots i_{m-2}} 
		+ \cdots 
		+ g_{i_{m-3} i_{m-2}}(f^{(m-2)}_{00})_{i_1 \dots i_{m-4}}
		\Big].
		\]
		For the other brackets, the same reasoning yields 
		\[
		\sum_{k=1}^{n} \binom{m}{2}^{-1} \Big[
		g_{i_1 i_2}(f^{(m-2)}_{kk})_{i_3 \dots i_{m-2}} 
		+ \cdots 
		+ g_{i_{m-3} i_{m-2}}(f^{(m-2)}_{kk})_{i_1 \dots i_{m-4}}
		\Big].
		\]
		Together, these two expressions contribute 
		\[
		\binom{m-2}{2} \binom{m}{2}^{-1} 
		\big[ \gc J f^{(m-2)} \big]_{i_1 \dots i_{m-2}}.
		\]
		
		Now we are left with the following terms:
		\[
		\begin{aligned}
			&-c^2 \binom{m}{2}^{-1} 
			\Big[ g_{i_1 0} f^{(m-2)}_{i_2 \dots i_{m-2} 0} 
			+ \cdots 
			+ g_{i_{m-2} 0} f^{(m-2)}_{i_1 \dots i_{m-3} 0} \Big] 
			+ \cdots \\ 
			&\quad + \sum_{k=1}^{n} \binom{m}{2}^{-1} 
			\Big[ g_{i_1 k} f^{(m-2)}_{i_2 \dots i_{m-2} k} 
			+ \cdots 
			+ g_{i_{m-2} k} f^{(m-2)}_{i_1 \dots i_{m-3} k} \Big].
		\end{aligned}
		\]
		
		Rewriting, we obtain
		\[
		\begin{aligned}
			&= \binom{m}{2}^{-1} \Big[ -c^2 g_{i_1 0} f^{(m-2)}_{i_2 \dots i_{m-2} 0} 
			+ \sum_{k=1}^n g_{i_1 k} f^{(m-2)}_{i_2 \dots i_{m-2} k} \Big] + \cdots \\
			&\quad + \binom{m}{2}^{-1} \Big[ -c^2 g_{i_{m-2} 0} f^{(m-2)}_{i_1 \dots i_{m-3} 0} 
			+ \sum_{k=1}^n g_{i_{m-2} k} f^{(m-2)}_{i_1 \dots i_{m-3} k} \Big].
		\end{aligned}
		\]
		By the definition of \(g\), in each bracket exactly one term survives, with coefficient \(+1\). Hence, the total contribution is
		\[
		(m-2)\binom{m}{2}^{-1} f^{(m-2)}_{i_1 \dots i_{m-2}}.
		\]
		
		Putting everything together, we obtain
		\[
		J \gc f^{(m-2)} 
		= \binom{m-2}{2}\binom{m}{2}^{-1} \gc J f^{(m-2)} 
		+ \binom{m}{2}^{-1} (m+n-1) \,f^{(m-2)}.
		\]
		This completes the proof of the first part.
		
		For the second part, we prove the claim by considering particular choices of indices.  
		First, let all indices vanish, i.e.\ \(i_1, \dots, i_{m-2} = 0\). Then  
		\[
		\big[ \gc f^{(m-2)} \big]_{0 \cdots 0} = 0 
		\;\;\implies\;\; f^{(m-2)}_{0 \cdots 0} = 0.
		\]
		
		Next, allow exactly one index to be nonzero. Without loss of generality, let \(i_1 \neq 0\) and all others vanish. Then  
		\[
		\big[ \gc f^{(m-2)} \big]_{i_1 0 \cdots 0} = 0 
		\;\;\implies\;\; f^{(m-2)}_{i_1 0 \cdots 0} = 0.
		\]
		By the symmetry of \(f^{(m-2)}\), every component with exactly one nonzero index must vanish. Proceeding inductively, increasing the number of nonzero indices at each step, the same argument shows that all corresponding components vanish. Hence, every component of \(f^{(m-2)}\) is zero, completing the proof.
	\end{proof}

\end{document}
